\scanpage{302}
\mspage{\textcircled{1}}
% slanted
\margin{1975 ou 1976}
\[
 \mathbb{\Delta}=\text{cat. des simplexes }\Delta^n\ (n\geq0)
 \qquad\substack{\Delta^n=[0,n]\subset\mathbb{N}\\\text{avec rel. d'ordre total}}
\]
\[
 \widehat{\mathbb{\Delta}}=\mathrm{Ss}
 =\underline{\mathrm{Hom}}(\mathbb{\Delta}^o,\mathrm{E}\mathrm{ns}))
 =\mathbb{E}\mathrm{ns}_*
\]

Plus généralement, si $A$ est une catégorie, on pose
\[
 \begin{aligned}
 A_*&=\underline{\mathrm{Hom}}(\mathbb{\Delta}^o,\mathbb{E}\mathrm{ns})\\
 A^*&=\underline{\mathrm{Hom}}(\mathbb{\Delta},\mathrm{Ens})
 \end{aligned}
\]
donc
\[
 (A^o)^*=(A_*)^o\qquad A^*\simeq((A^o)_*)^o
\]
où l'exposant $o$ désigne le \inserted{passage à la} cat. opposée.

Les objets de $A_*$ (resp. $A^*$) sont considérés comme des structures algébriques de type simple (sur une infinité d'objets de base) dans $A$, les ``objets $1/2$ simpliciaux'' resp ``$1/2$ cosimpliciaux''. \inserted{On} écrira $\mathbb{E}\mathrm{ns}_*$ quand on a en vue cet aspect, et $\widehat{\mathbb{\Delta}}$ \inserted{ou $\mathrm{Ss}$} quand on a plutôt le pt de vue ``objet d'un topos'' ou ``faisceaux sur un topos''\deleted{\unsure{,}}\textsuperscript{$\circ\!/$}% 
% margin note written sideways
\margin{\textsuperscript{$\circ\!/$} La notation $\widehat{\mathbb{\Delta}}$ peut être utilisée de préf. quand on veut garder à l'esprit la nature algébrique très particulière de ce topos, la notation $\mathrm{Ss}$ quand \deleted{\unsure{on}} veut l'oublier au profit de sa signification topologique.}
Tous les développements qui suivent ont pour objet d'étudier ces objets (l'équivalent combinatoire des espaces topologiques) et leurs ``types d'homotopie'', via des invariants qu'on peut leur associer, qui en dépendent de façon covariante ou contravariante. Nous avons ici en vue une étude systématique, dans l'esprit de l'algèbre universelle, de ces invariants, (qui sont donc des foncteurs sur $\widehat{\mathbb{\Delta}}\unsure{\simeq|=}\mathbb{E}\mathrm{ns}_*$) \deleted{et des} \inserted{de leurs structures et des} opérations qui peuvent être définies sur elles, permettant d'en déduire certains complexes \unsure{:} à partir d'autres plus simples

% archive sheet number: 303
